\begin{table}[!htbp]
\centering
\caption{Comparação das seleções, distâncias e predições entre modelos.}
\label{tab:res-contrastes}
\begin{singlespace}
\normalsize
\setlength{\tabcolsep}{5pt}
\renewcommand{\arraystretch}{1.15}
\begin{tabular}{lrrrrr}
\hline\hline
\textbf{Contraste} & \textbf{\shortstack{Casos\\comuns}} & \textbf{\shortstack{Atributos\\comuns}} & \textbf{Jaccard} & \textbf{$\rho$ distâncias} & \textbf{\shortstack{Concordância\\(\%)}} \\
\hline
55--74: controle & 55 & 4 & 0,250 & 0,8383 & 96,97 \\
55--74: topologia & 55 & 2 & 0,111 & 0,9540 & 99,09 \\
74: controle--topologia & 74 & 0 & 0,000 & 0,8484 & 94,59 \\
55: controle--topologia & 55 & 4 & 0,250 & 0,9018 & 98,48 \\
\hline\hline
\end{tabular}
\par\smallskip
{\small Nota: o Jaccard compara os conjuntos de dez atributos. A correlação de Spearman compara distâncias dos mesmos identificadores, com quatro casas decimais. A concordância compara predições entre modelos e não os acertos contra os rótulos.\par}
{\fonte{elaboração própria a partir dos dados experimentais.}}
\end{singlespace}
\end{table}
